\section{Arduinoのスケッチ}
作成したArduinoのスケッチを以下に示す．

\vspace{1cm}\begin{lstlisting}[caption=sound\_tx.ino]
const int maxmoji = 256 * 4;  // 最大文字数256文字分（ASCII）*4桁
int base4[maxmoji];           // 4進数の結果を格納する配列
int resultIndex = 0;
int time = 5;

void setup() {
  Serial.begin(9600);
  pinMode(A0, OUTPUT);
  pinMode(A1, OUTPUT);
  pinMode(A2, OUTPUT);
  pinMode(A3, OUTPUT);
  pinMode(A4, OUTPUT);
}

void loop() {
  if (Serial.available() > 0) {
    String input = Serial.readStringUntil('\n');
    resultIndex = 0;

    for (int i = 0; i < input.length(); i++) {

      char c = input.charAt(i);
      byte ascii = (byte)c;
      getBase4(ascii);
    }

    Base4send();
    
    for (int i = 0; i < resultIndex; i++) {
      Serial.print(base4[i]);
    }
    Serial.println("");
  }
}

void getBase4(byte b) {
  int amari[4];
  for (int i = 3; i >= 0; i--) {
    amari[i] = b % 4;
    b = b / 4;
  }
  for (int i = 0; i < 4; i++) {
    base4[resultIndex++] = amari[i];
  }
}

void Base4send() {
  digitalWrite(A4, HIGH);

  for (int k = 0; k < resultIndex; k++) {
    int pin;
    if (base4[k] == 0) {
      pin = A0;
    } else if (base4[k] == 1) {
      pin = A1;
    } else if (base4[k] == 2) {
      pin = A2;
    } else {
      pin = A3;
    }

    digitalWrite(pin, HIGH);
    delay(time);              // この時間内に受信側が読み取り
    digitalWrite(pin, LOW);
  }

  digitalWrite(A4, LOW);    // 送信終了通知
}

\end{lstlisting}
% -----------------------------------------------------------

\vspace{1cm}\begin{lstlisting}[caption=sound\_tx\_itg.ino]
#include"hack.h"

const int maxmoji = 256 * 4;  // 最大文字数256文字分（ASCII）*4桁
int base4[maxmoji];           // 4進数の結果を格納する配列
int resultIndex = 0;          // 4進数変換したあとの文字数+1

// 音送信用の変数
const int speaker_pin = 3;  //スピーカー接続ピン
unsigned long t = 300;  // 音を鳴らす時間(ms)
unsigned long GI = t * 0.25; // GI長(tの25%)
unsigned int freq[16] = {2000, 2100, 2200, 2300, 2400, 2500, 2600, 2700, 2800, 2900, 3000, 3100, 3200, 3300, 3400, 3500};//鳴らす周波数を設定（最低周波数を開始ビットに）
unsigned int sfreq = 3600;
String str[16] = {"00","01","02","03","10","11","12","13","20","21","22","23","30","31","32","33"} ; // 取得した4進数2桁
double txtime = 0;

// 光受信用の変数
const int detectPin = A4; // A4の光（送信中）を検出
bool inReceiving = false; // 現在受信中かどうか
const int read_time = 10; // 読み取る間隔[ms]

void setup() {
  Serial.begin(9600);
  pinMode(speaker_pin, OUTPUT);
  for (int i = 0; i < 4; i++) {
    pinMode(dataPins[i], INPUT);
  }
}

void loop() {
  int detectValue = analogRead(detectPin);
  if (detectValue > threshold) {
    // 送信中
    if (!inReceiving) {
      resultIndex = 0;      // 初回のみ初期化
      inReceiving = true;
    }

    int val = detectBitByDuration(read_time);  // msごとの読み取り
    if (val != -1) {
      base4[resultIndex++] = val; // 桁数を数える
    }
  }else{
    // 送信終了時
    if (inReceiving) {
      txtime = 0;
      txtime = millis();
      String s = "";

      // 開始用の周波数を出力
      syncTone(speaker_pin, sfreq, t, GI);
      syncTone(speaker_pin, sfreq, t, GI);

      for (int i=0; i<resultIndex; i++){
        s += String(base4[i]);
      }
      for (int k=0; k<resultIndex; k+=2){
        // 前から2桁ずつ音を鳴らす
        for (int l=0; l<16; l++){
          if(s.substring(k, k+2) == str[l]){
            tone(speaker_pin, freq[l]);
            delay(t);
            noTone(speaker_pin);
            delay(GI);
          }
        }
      }
      delay(GI);

      // 終了用の周波数を出力
      syncTone(speaker_pin, sfreq, t, GI);
      syncTone(speaker_pin, sfreq, t, GI);

      // 確認用シリアル表示（読み取り）
      Serial.print("受信データ: ");
      for (int m=0; m<resultIndex; m++) {
        Serial.print(base4[m]);
      }
      Serial.println("");
      Serial.println("-----------------");
      Serial.print("Base4: ");
      for (int i = 0; i < resultIndex; i++) {
        Serial.print(base4[i]); // 4進数の左から表示
        s += String(base4[i]);
      }
      Serial.println("");
      for (int l=0; l<resultIndex; l+=2){
        for (int m=0; m<16; m++){
          if(s.substring(l, l+2) == str[m]){
            String foo = s.substring(l, l+2);
            int bar = ((foo.charAt(0) - '0') * 4 + (foo.charAt(1) - '0'));
            Serial.print(freq[bar]);
            Serial.println(" Hz");
          }
        }
      }
      Serial.println("------------------------");
      inReceiving = false;  // フラグを戻す
    }
  }
}
\end{lstlisting}

%--------------------------------------------------

\vspace{1cm}\begin{lstlisting}[caption=hack.h]
    #ifndef HACK_H
#define HACK_H

#include <Arduino.h>

const int dataPins[4] = {A0, A1, A2, A3};
const int threshold = 50;

int detectBitByDuration(int durationMs);
void syncTone(int pin, unsigned int  freqency, unsigned long delay, unsigned long GI);

#endif
\end{lstlisting}
%-----------------------------------------------

\vspace{1cm}\begin{lstlisting}[caption=hack.cpp]
 #include "hack.h"

int detectBitByDuration(int durationMs) {
  int lightDetected[4] = {0, 0, 0, 0};
  unsigned long start = millis();

  while (millis() - start < durationMs) {
    for (int i = 0; i < 4; i++) {
      if (analogRead(dataPins[i]) > threshold) {  // 光を検出したらカウント
        lightDetected[i]++;
      }
    }
  }

  // 最も多く光を検出したピンを選ぶ（ノイズ対策）
  int maxIndex = -1;
  int maxVal = 0;
  for (int i = 0; i < 4; i++) {
    if (lightDetected[i] > maxVal) {
      maxVal = lightDetected[i];
      maxIndex = i;
    }
  }

  // しきい値以上の光検出があれば有効
  if (maxVal > 3) {
    return maxIndex;
  } else {
    return -1;  // ノイズとみなして無視
  }
}

void syncTone(int pin, unsigned int  freqency, unsigned long delayt, unsigned long GI){
    tone(pin, freqency);
    delay(delayt);
    noTone(pin);
    delay(GI);
}

\end{lstlisting}
%----------------------------------------------------------

\vspace{1cm}\begin{lstlisting}[caption=vibration\_tx\_itg.ino]
#include <arduinoFFT.h>

// --- 受信側（マイク）設定 ---
const uint16_t samples = 256;              // FFTサンプル数
const double samplingFrequency = 8000.0;   // サンプリング周波数 (Hz)
const double samplingIntervalUs = 1000000.0 / samplingFrequency; // サンプリング間隔 (マイクロ秒)
const uint8_t micPin = A0;                 // マイク入力用アナログピン

double vReal[samples];
double vImag[samples];

ArduinoFFT FFT = ArduinoFFT(vReal, vImag, samples, samplingFrequency);

String decodedData = "";
bool receiving = false;
int startSignalCount = 0;
int endSignalCount = 0;

// --- 受信側データマッピング ---
String str[16] = {"00","01","02","03","10","11","12","13","20","21","22","23","30","31","32","33"};
unsigned int binaryarray[] = {0, 1, 10, 11, 100, 101, 110, 111, 1000, 1001, 1010, 1011, 1100, 1101, 1110, 1111};
const int maxlit = 256 * 4; // 受信データ格納用配列の最大サイズ
int binary[maxlit]; // パディングされた2進数を整数として格納（主にデバッグ表示用）
int ReceivedLength;

// --- 送信側（振動モーター）設定 ---
const int motorPin = 5; // 振動モーター用デジタルピンをD5に設定
const unsigned long bitDuration = 200;  // 1ビットの持続時間 (ms)
const int signalBitCount = 7;           // 送信開始信号ビット数（1を7回）
const int endSignalBitCount = 10;       // 送信終了信号ビット数（0を10回）

// --- 関数プロトタイプ宣言 ---
// 受信側
double detectFrequency();
int decodeFrequency(double freq);
// 受信した4進数文字列をASCIIバイト列に変換し、各バイトを振動送信する関数
void convertAndSendVibration(String quadStr);

// 送信側
void sendVibrationBit(int bitValue);
void sendByte(byte data);
void sendStartSignal();
void sendEndSignal();


void setup() {
  Serial.begin(9600);           // シリアル通信の初期化
  pinMode(motorPin, OUTPUT);    // モーターピンを出力に設定
  digitalWrite(motorPin, LOW);  // モーターを初期状態はOFFにする
}

void loop() {
  // --- 受信処理ブロック ---
  double freq = detectFrequency();
  int decoded = decodeFrequency(freq);
  int binaryIndex = 0; // binary配列の書き込みインデックスをここで宣言・初期化
  for (int i = 0; i < maxlit; i++) {
    binary[i] = 0; // binary配列を毎回リセット
  }

  if (decoded == -1) {  // 開始/終了シグナル
    if (!receiving) {
      // 通信開始前
      startSignalCount++;
      if (startSignalCount >= 2) { // 2回以上受け取ったら開始
        decodedData = "";  // 開始時に初期化
        receiving = true;
        startSignalCount = 0;
        endSignalCount = 0;
        Serial.println("----- Start receiving -----");
      } else {
        Serial.print("Start signal received ");
        Serial.print(startSignalCount);
        Serial.println(" time(s). Waiting for 2nd signal...");
      }
    } else {
      // 通信中
      endSignalCount++;
      if (endSignalCount >= 2) { // 2回以上受け取ったら終了
        receiving = false;
        startSignalCount = 0;
        endSignalCount = 0;
        Serial.println("----- End receiving -----");
        Serial.print("Received(row): ");
        Serial.println(decodedData);

        Serial.print("Decoded ASCII: "); // 新しいラベル
        String fullBinaryStringForDisplay = "";
        int currentReceivedLengthForDisplay = decodedData.length();

        for(int i = 0; i < currentReceivedLengthForDisplay; i += 2){
          String currentQuad = decodedData.substring(i, i + 2);
          bool found = false;

          for(int j = 0; j < 16; j++){
            if(currentQuad == str[j]){
              String binStr = String(binaryarray[j]);
              while (binStr.length() < 4) {
                  binStr = "0" + binStr;
              }
              fullBinaryStringForDisplay += binStr;
              found = true;
              break;
            }
          }
        }

        // 8ビットごとにASCII文字に変換して表示
        for (int i = 0; i + 7 < fullBinaryStringForDisplay.length(); i += 8) {
          String byteBinaryStr = fullBinaryStringForDisplay.substring(i, i + 8);
          long byteValue = 0;
          for (int k = 0; k < 8; k++) {
            if (byteBinaryStr.charAt(k) == '1') {
              byteValue += (1 << (7 - k));
            }
          }
          Serial.print((char)byteValue); // 文字を直接表示
        }
        Serial.println(); // デコードされた文字の表示後に改行

        // 受信した4進数データをASCIIに変換し、そのバイナリ表現を振動として送信
        convertAndSendVibration(decodedData);
        
        // --- 以下はデバッグ表示用（受信機のデバッグ出力として残す） ---
        ReceivedLength = decodedData.length();
        Serial.print("Received length(row):  ");
        Serial.println(ReceivedLength);

        for(int i=0; i<ReceivedLength; i+=2){
          String currentQuad = decodedData.substring(i, i+2);
          for(int j=0; j<16; j++){
            if(currentQuad == str[j]){
              if(binaryIndex < (sizeof(binary)/sizeof(binary[0]))){
                String paddedBinaryStr = String(binaryarray[j]);
                while (paddedBinaryStr.length() < 4) {
                    paddedBinaryStr = "0" + paddedBinaryStr;
                }
                binary[binaryIndex] = paddedBinaryStr.toInt();
                Serial.print("binaryarray (padded int): ");
                Serial.println(binary[binaryIndex]);
                binaryIndex++;
              }else{
                Serial.println("Error: Storage overflow for binary array.");
              }
              for(int l=0; l<binaryIndex; l++){
                Serial.print(binary[l]);
              }
              Serial.println("");
              Serial.print("binary (original value): ");
              Serial.println(binaryarray[j]);
              break;
            }
          }
        }
        Serial.print("finally binary[] contents (padded int): [");
        for (int i = 0; i < binaryIndex; i++) {
          Serial.print(binary[i]);
          if (i < binaryIndex - 1) {
            Serial.print(", ");
          }
        }
        Serial.println("]");
        // -----------------------------------------------------------------

      } else {
        Serial.print("End signal reveived ");
        Serial.print(endSignalCount);
        Serial.println(" time(s). Waiting for 2nd signal to end...");
      }
    }
  } else if (receiving && decoded >= 0 && decoded <= 15) {
    // データ受信中
    endSignalCount = 0; // データが来たら終了カウントをリセット
    char quadDigit1 = '0'+(decoded/4);
    char quadDigit2 = '0'+(decoded%4);
    String quadStr = String(quadDigit1) + String(quadDigit2);
    Serial.print("Adding quad: ");
    Serial.println(quadStr);
    decodedData += quadStr;
  } else if (!receiving && decoded >= 0 && decoded <= 15) {
    // 通信開始前にデータが来ても無視し、開始カウントをリセット
    startSignalCount = 0;
  }
  delay(290); // 次の音まで待機
}

// --- 受信側関数定義 ---

// FFTでピーク周波数を検出する関数
double detectFrequency() {
  unsigned long targetTime = micros();
  for (int i = 0; i < samples; i++) {
    while (micros() < targetTime) {}
    vReal[i] = analogRead(micPin);
    vImag[i] = 0;
    targetTime += samplingIntervalUs;
  }

  // DCオフセット除去
  double mean = 0;
  for (int i = 0; i < samples; i++) {
    mean += vReal[i];
  }
  mean /= samples;
  for (int i = 0; i < samples; i++) {
    vReal[i] -= mean;
  }

  FFT.windowing(FFT_WIN_TYP_HAMMING, FFT_FORWARD);
  FFT.compute(FFT_FORWARD);
  FFT.complexToMagnitude();

  double peak = 0;
  int peakIndex = 0;
  for (int i = 1; i < samples / 2; i++) {
    if (vReal[i] > peak) {
      peak = vReal[i];
      peakIndex = i;
    }
  }
  return (peakIndex * samplingFrequency) / samples;
}

// 周波数から4進数2桁データをデコード（0〜15）
int decodeFrequency(double freq) {
  int freqInt = (int)(freq + 0.5);

  if (freqInt >= 3570 && freqInt <= 3630) return -1; // 開始/終了信号

  for (int i = 0; i < 16; i++) {
    int expected = 2000 + 100 * i;
    if (freqInt >= expected - 49 && freqInt <= expected + 49) {
      return i;
    }
  }
  return -2; // 無効な周波数
}

// 受信した4進数文字列をASCIIバイト列に変換し、各バイトを振動送信する関数
void convertAndSendVibration(String quadStr) {
  String fullBinaryString = "";
  int currentReceivedLength = quadStr.length();

  // 4進数チャンクをパディングされた4ビットの2進数文字列に変換して連結
  for(int i = 0; i < currentReceivedLength; i += 2){
    String currentQuad = quadStr.substring(i, i + 2);
    bool found = false;

    for(int j = 0; j < 16; j++){
      if(currentQuad == str[j]){
        String binStr = String(binaryarray[j]);
        while (binStr.length() < 4) {
            binStr = "0" + binStr;
        }
        fullBinaryString += binStr;
        found = true;
        break;
      }
    }
    if (!found) {
      Serial.print("警告: 未知の4進数チャンク '");
      Serial.print(currentQuad);
      Serial.println("' が検出されました。スキップして '0000' を挿入します。");
      fullBinaryString += "0000"; // 未知のチャンクの場合、デフォルトで0000を挿入
    }
  }

  Serial.print(" Binary string (for vibration): ");
  Serial.println(fullBinaryString);
  Serial.print(" Converted ASCII (for vibration): ");
  Serial.println("");

  // まず送信開始信号を送る
  sendStartSignal();
  delay(bitDuration); // 送信開始信号後の待機

  // 8ビットごとにASCII文字に変換し、そのバイトを振動送信
  for (int i = 0; i + 7 < fullBinaryString.length(); i += 8) {
    String byteBinaryStr = fullBinaryString.substring(i, i + 8);
    long byteValue = 0;

    for (int k = 0; k < 8; k++) {
      if (byteBinaryStr.charAt(k) == '1') {
        byteValue += (1 << (7 - k));
      }
    }
    Serial.print((char)byteValue); // シリアルモニターに文字を表示 (ASCII)

    // デコードされたバイトを振動として送信
    sendByte((byte)byteValue);
  }
  Serial.println();

  // 全てのバイト送信後に送信終了信号を送る
  sendEndSignal();
  Serial.println("Vibration transmission completed");
  digitalWrite(motorPin, LOW); // モーターをOFFにする
}


// --- 送信側関数定義 ---

// 1ビットの値を振動として出力する関数
void sendVibrationBit(int bitValue) {
  digitalWrite(motorPin, bitValue == 1 ? HIGH : LOW);
  delay(bitDuration);
  digitalWrite(motorPin, LOW);  // 振動のOFFを明示
}

// 1バイトのデータをビット列に分解して振動として送信する関数
void sendByte(byte data) {
  Serial.print("  Bits for '");
  Serial.print((char)(data & 0x7F)); // 送信する文字（7ビットASCII）を表示
  Serial.print("': ");
  for (int bit = 6; bit >= 0; bit--) { // 7ビットASCIIの範囲でループ (MSBからLSBへ)
    int bitValue = (data >> bit) & 0x01; // ビット値を取得
    sendVibrationBit(bitValue);          // 振動として送信
    Serial.print(bitValue);              // シリアルモニターにビット値を表示
  }
  Serial.println();
}

// 送信開始信号を送る関数
void sendStartSignal() {
  Serial.println("--- 振動送信開始信号 ---"); // シリアルモニターに表示
  for (int i = 0; i < signalBitCount; i++) {
    sendVibrationBit(1); // 1 (HIGH) を signalBitCount 回送信
  }
}

// 送信終了信号を送る関数
void sendEndSignal() {
  Serial.println("--- 振動送信終了信号 ---"); // シリアルモニターに表示
  for (int i = 0; i < endSignalBitCount; i++) {
    sendVibrationBit(0); // 0 (LOW) を endSignalBitCount 回送信
  }
}
\end{lstlisting}
%------------------------------------

\vspace{1cm}\begin{lstlisting}[caption=vibration\_rcv.ino]

const int xPin = A0;
const int yPin = A1;
const int zPin = A2;

const unsigned long bitDuration = 230;           // 1ビットの持続時間（ms）
const float vibrationThreshold = 12.9;           // 振動検出のしきい値
const int signalBitCount = 5;                    // 開始信号（1*7回）
const int endSignalBitCount = 11;                // 終了信号（0*10回）
const int maxBits = 512;                         // 最大ビット数バッファ

// 加速度センサから3軸データを読み取る
void readXYZ(int &x, int &y, int &z) {
  x = analogRead(xPin);
  y = analogRead(yPin);
  z = analogRead(zPin);
}

// 指定時間内に振動があったかを検出
bool detectVibration(unsigned long duration) {
  unsigned long start = millis();
  int baseX, baseY, baseZ;
  readXYZ(baseX, baseY, baseZ);

  int vibrationCount = 0;
  int samples = 0;

  while (millis() - start < duration) {
    int xNow, yNow, zNow;
    readXYZ(xNow, yNow, zNow);
    float delta = sqrt(
      pow(xNow - baseX, 2) +
      pow(yNow - baseY, 2) +
      pow(zNow - baseZ, 2)
    );
    if (delta > vibrationThreshold) vibrationCount++;
    samples++;
    delay(5);  // 過剰サンプリング防止
  }

  return (vibrationCount > samples / 2);
}

// 開始信号（1*7）を検出
bool detectStartSignal() {
  Serial.println("Checking start signal...");
  for (int i = 0; i < signalBitCount; i++) {
    if (!detectVibration(bitDuration)) {
      Serial.println("Start signal failed");
      return false;
    }
    Serial.print("1");
  }
  Serial.println("\n[Start Detected]");
  return true;
}

// 終了信号（最後の10ビットが0）を検出
bool detectEndSignal(const bool *bits, int len) {
  if (len < endSignalBitCount) return false;
  for (int i = len - endSignalBitCount; i < len; i++) {
    if (bits[i]) return false;
  }
  return true;
}

// --- 初期化 ---
void setup() {
  Serial.begin(9600);
  pinMode(xPin, INPUT);
  pinMode(yPin, INPUT);
  pinMode(zPin, INPUT);
  Serial.println("Ready to receive vibration-based communication");
}

// --- メインループ ---
void loop() {
  static bool bitsBuffer[maxBits];
  static int bitsCount = 0;

  Serial.println("Waiting for start signal...");
  while (!detectStartSignal());  // 開始信号を待機

  delay(3*bitDuration);

  bitsCount = 0;
  unsigned long bitStart;

  // --- ビット受信ループ ---
  while (true) {
    bitStart = millis();
    bool bit = detectVibration(bitDuration);
    bitsBuffer[bitsCount++] = bit;
    Serial.print(bit ? '1' : '0');

    // 終了信号検出
    if (bitsCount >= endSignalBitCount && detectEndSignal(bitsBuffer, bitsCount)) {
      Serial.println("\n[End Detected]");
      break;
    }

    // ビット数上限
    if (bitsCount >= maxBits) {
      Serial.println("\n[Error] Buffer overflow. Resetting.");
      return;
    }

    // 次のビットまで待機（正確な周期維持）
    while (millis() - bitStart < bitDuration) {
      delay(1);
    }
  }

  // --- 7ビットASCIIに復元 ---
  int dataBits = bitsCount - endSignalBitCount;
  int usableBits = (dataBits / 7) * 7;

  Serial.print("Decoded message: ");
  for (int i = 0; i + 6 < usableBits; i += 7) {
    byte c = 0;
    for (int b = 0; b < 7; b++) {
      if (bitsBuffer[i + b]) {
        c |= (1 << (6 - b));  // MSBからLSB順に復元
      }
    }
    Serial.print((char)c);
  }
  Serial.println();
}

\end{lstlisting}
